% -*- TeX-PDF-mode: t; IspellDict: english -*-
\documentclass[aspectratio=169]{beamer}
%
\newcommand{\figPath}{figures}
\newcommand{\imgPath}{images}
%
% Core ------------------------------------------------------------------------%
\usetheme{Lumen}
\usepackage[spanish,english]{babel}

% Math (requires being first) -------------------------------------------------%
% \usepackage{amsmath}
% \usepackage{amsthm}
\usepackage{galois}
% \usepackage{marvosym}
\usepackage[ruled,vlined]{algorithm2e}
\usepackage{empheq}

% Layout & Document -----------------------------------------------------------%
\usepackage{transparent}
\usepackage{graphicx}
\usepackage{xspace}
\usepackage[dvipsnames]{xcolor}
%
\definecolor[named]{MyPurple}{cmyk}{0.55,1,0,0.15}
\definecolor[named]{MyDarkBlue}{cmyk}{1,0.58,0,0.21}
\definecolor[named]{LightGray}{gray}{0.9}

% Fonts -----------------------------------------------------------------------%
\usepackage[T1]{fontenc}
\usepackage[utf8]{inputenc}
\usepackage[scaled=.97,zerostyle=b]{newtxtt}
\usepackage[scr=boondox]{mathalpha}
%
\font\bboldfonttwelve=bbold12 at 12pt%
\font\bboldfontten=bbold12 at 10pt%
\font\bboldfonteight=bbold12 at 8pt%
\renewcommand{\mathbb}[1]{%
  \mathchoice%
    {\mbox{\bboldfontten#1}}%
    {\mbox{\bboldfontten#1}}%
    {\mbox{\scriptsize\bboldfontten#1}}%
    {\mbox{\tiny\bboldfonteight#1}}%
}
%
\usepackage{amsbsy}

% Listings --------------------------------------------------------------------%
\usepackage{xparse}
\usepackage{listings}

% Tables ----------------------------------------------------------------------%
\usepackage{booktabs}
\usepackage{multirow}

% Hyperlinks & References -----------------------------------------------------%
\usepackage{url}
%
\hypersetup{%
  colorlinks=true,%
  linkcolor=black,%
  citecolor=MyDarkBlue,%
  urlcolor=MyDarkBlue,%
  filecolor=MyDarkBlue%
}

% % Beamer ----------------------------------------------------------------------%
\usepackage{appendixnumberbeamer}

% Wants to be Last? -----------------------------------------------------------%
\usepackage{csquotes}

%%% Local Variables:
%%% mode: LaTeX
%%% TeX-master: "main"
%%% End:

% (included file)

%
\newcommand{\inter}{\ensuremath{{\pmb{\mathfrak{I}}}}}
\newcommand{\supInter}[1]{\ensuremath{#1}}
%
\newcommand{\sem}[1]{\ensuremath{\mathbb{[}\,#1\,\mathbb{]}}}
\newcommand{\asem}[1]{\ensuremath{\pmb{\mathscr{A}}\sem{#1}}}
\newcommand{\bsem}[1]{\ensuremath{\pmb{\mathscr{B}}\sem{#1}}}
\newcommand{\esem}[1]{\ensuremath{\pmb{\mathscr{E}}\sem{#1}}}
\newcommand{\ssem}[1]{\ensuremath{\pmb{\mathscr{S}}\sem{#1}}}


%%% Local Variables:
%%% mode: LaTeX
%%% TeX-master: "main"
%%% End:

%
\title{Formal Semantics and its Approximation}
\subtitle{An Introduction}
%
\author{\centering Marco Ciccal\`{e}\\{\fontfamily{cmtt}\selectfont\scriptsize marco.ciccale@imdea.org}}
\date{\textbf{Aut.~Coffee} -- March 10, 2026}
\institute{%
  Universidad Politécnica de Madrid (UPM), Madrid Spain\\
  IMDEA Software Institute, Madrid, Spain
}
%
\begin{document}
%
\begin{frame}[plain]
  \titlepage
\end{frame}
%
\addtocounter{framenumber}{-1}
%
\section{Programs}

\begin{frame}{Syntax}
  \begin{itemize}
  \item We will work with a simple imperative programming language.
  \item We start by defining how programs must be \emph{written}.
  \end{itemize}
  %
  \pause
  %
  \begin{block}{Syntax of expressions}
    \vspace*{-2em}
    \begin{flalign*}
      \texttt{x}, \texttt{y}, \ldots \quad \in \quad \mathbb{V} && \text{variables}\\
      \texttt{E} \quad \in \quad \mathbb{E} & \quad ::= \quad \texttt{1}~
                                              |~\texttt{x}~
                                              |~\texttt{E}_1 \texttt{ - } \texttt{E}_2
                                            & \text{arithmentic expressions}\\
                                            & \quad \quad| \;\;\:\quad \texttt{tt}~
                                              |~\texttt{ff}~
                                              |~\texttt{E}_1 \texttt{ < } \texttt{E}_2~
                                              |~\texttt{E}_1 \texttt{ nand } \texttt{E}_2
                                            & \text{Boolean expressions}
    \end{flalign*}
  \end{block}
  %
  \vspace{-1em}\pause
  %
  \begin{block}{Syntax of statements}
    \vspace*{-2em}
    \begin{align*}
        \texttt{S} \quad \in \quad \mathbb{S} & \quad
        ::= \quad \texttt{skip}~
          |~\texttt{x := E}~
          |~\texttt{S}_1 \texttt{; } \texttt{S}_2~
          |~\texttt{\textbf{if} (E) \textbf{then }} \texttt{S}_1 \texttt{\bfseries~else } \texttt{S}_2~
          |~\texttt{\textbf{while} (E) \textbf{do }} \texttt{S}
    \end{align*}
    \vspace*{-2em}
  \end{block}
  %
  \begin{itemize}
  \pause\item Purely \emph{syntactic} object, it has no \emph{meaning}
    (\emph{i.e.}, no semantics) yet.
  \pause\item Defined \emph{compositionally}, \emph{i.e.}, composite
    elements such as \texttt{E}\textsubscript{1}\texttt{ - }\texttt{E}\textsubscript{2}
    have a unique decomposition into \emph{strictly smaller} elements.
  \end{itemize}
\end{frame}
%
\begin{frame}{Semantics}
  \begin{itemize}
  \item We now turn our heads to the meaning of programs: \emph{semantics}.
  \item Semantics can come in many different flavors, and they complement each other.
  \end{itemize}
  %
  \pause\medskip
  %
  {\small\begin{minipage}[t]{.32\textwidth}
    \begin{center}
      \textbf{Operational\\\cite{plotkin-sopsem}}
    \end{center}
    %
    \vspace{-.8em}
    %
    \begin{itemize}
    \item Defines \emph{how} a program is executed.
    \item \emph{Small-} or \emph{big-step} definitions.
    \end{itemize}
  \end{minipage}
  %
  \hfill
  %
  \begin{minipage}[t]{.32\textwidth}
    \begin{center}
      \textbf{Denotational\\\cite{scott71}}
    \end{center}
    %
    \vspace{-.8em}
    %
    \begin{itemize}
    \item Models the \emph{result} of executing a program.
    \item \emph{Compositional} definitions.
    \end{itemize}
  \end{minipage}
  %
  \hfill
  %
  \begin{minipage}[t]{.32\textwidth}
    \begin{center}
      \textbf{Axiomatic\\\cite{hoare69}}
    \end{center}
    %
    \vspace{-.8em}
    %
    \begin{itemize}
    \item Allows to prove (partial) correctess of a program.
    \item Assertion-based definitions.
    \end{itemize}
  \end{minipage}}
  %
  \pause\bigskip
  %
  \begin{itemize}
  \item Herein, we concentrate in \emph{denotational} semantics, but
    the approaches can be adapted to any semantics flavor.
  \end{itemize}
  %
  % \item As \emph{numerals} (\emph{i.e.}, expressions) and
  %   \emph{numbers} (\emph{i.e.,} mathematical, abstract objects)
  %   differ, there is a distinction between the program itself and what
  %   it \emph{denotes}.
  % \end{itemize}
\end{frame}
%
\begin{frame}{Denotational semantics: expressions}
  \begin{itemize}
  \item We consider the \emph{structure}
    $\inter \triangleq ( \mathbb{Z}, \mathbb{B}; \supInter{1}, \supInter{\mathfrak{tt}}, \supInter{\mathfrak{ff}},
    \supInter{-}, \supInter{<}, \supInter{\uparrow})$ with the usual interpretations.
  \end{itemize}
  %
  \begin{block}{Denotational semantics of expressions}
    $\pmb{\mathscr{E}} : \mathbb{V} \to (\texttt{E} \in \mathbb{E})
    \to (\mathbb{V} \to \mathbb{Z} \cup \mathbb{B}) \rightharpoonup
    \mathbb{Z} \cup \mathbb{B}
    %
    \text{\small~~~~~where } \mathbb{vars}(\texttt{E}) \subseteq
    \mathbb{V}$\footnote{%
      We use a dependent type to express that we require the
      expression \texttt{E} to have all of its variables considered in the
      (implicit) set of variables $\mathbb{V}$.  A more detailed discussion on
      this matter can be found in~\cite{Cousot21-book}.% 
    }
    %
    \pause\begin{flalign*}
    \hspace*{-.5em}\begin{array}{lllll}
      \esem{\texttt{k}} \sigma
      &\triangleq
      &k
      &\in \mathbb{Z} \cup \mathbb{B}
      &\\[.5em]
      %
      \esem{\texttt{x}} \sigma
      &\triangleq
      &\sigma(\texttt{x})
      &\in \mathbb{Z} \cup \mathbb{B}
      &\\[.5em]
      %
      \esem{\texttt{E}_1 \texttt{ - } \texttt{E}_2} \sigma
      &\triangleq
      &v_1 \supInter{-} v_2
      &\in \mathbb{Z}
      &\text{%
        if $v_1 = \esem{\texttt{E}_1}\sigma \in \mathbb{Z}$,
           $v_2 = \esem{\texttt{E}_2}\sigma \in \mathbb{Z}$
      }\\[.5em]
      %
      \esem{\texttt{E}_1 \texttt{ < } \texttt{E}_2} \sigma
      &\triangleq
      &v_1 \supInter{<} v_2
      &\in \mathbb{B}
      &\text{%
        if $v_1 = \esem{\texttt{E}_1}\sigma \in \mathbb{Z}$,
           $v_2 = \esem{\texttt{E}_2}\sigma \in \mathbb{Z}$
      }\\[.5em]
      %
      \esem{\texttt{E}_1 \texttt{ nand } \texttt{E}_2} \sigma
      &\triangleq
      &v_1 \supInter{\uparrow} v_2
      &\in \mathbb{B}
      &\text{%
       if $v_1 = \esem{\texttt{E}_1}\sigma \in \mathbb{E}$,
           $v_2 = \esem{\texttt{E}_2}\sigma \in \mathbb{E}$
      }\\[.5em]
      %
      undefined &&&& otherwise
    \end{array}&&
    \end{flalign*}
  \end{block}
\end{frame}
%
%
\begin{frame}{Denotational semantics: statements}
  \begin{itemize}
  \item Now, we turn our heads to statements.
  \end{itemize}
  %
  \begin{block}{Denotational semantics of statements}
    $\pmb{\mathscr{S}} : \mathbb{V} \to (\texttt{S} \in \mathbb{S})
    \to (\mathbb{V} \to \mathbb{Z} \cup \mathbb{B}) \rightharpoonup
    (\mathbb{V} \to \mathbb{Z} \cup \mathbb{B})
    %
    \text{\small~~~~~where } \mathbb{vars}(\texttt{E}) \subseteq
    \mathbb{V}$
    %
    \pause\begin{flalign*}
    \hspace*{-.5em}\begin{array}{llll}
      \ssem{\texttt{skip}} \sigma
      &\triangleq
      &\sigma
      &\\[.5em]
      %
      \ssem{\texttt{x := E}} \sigma
      &\triangleq
      &\sigma[v/\texttt{x}]
      &\text{%
        if $v = \esem{\texttt{E}}\sigma$
      }\\[.5em]
      %
      \ssem{\texttt{S}_1 \texttt{; } \texttt{S}_2}
      &\triangleq
      &\ssem{\texttt{S}_2} \circ \ssem{\texttt{S}_1}
      &\\[2em]
      %
      \ssem{\texttt{\textbf{if} (E) \textbf{then }} \texttt{S}_1 \texttt{\bfseries~else } \texttt{S}_2} \sigma
      &\triangleq
      &\begin{cases}
        \ssem{\texttt{S}_1} \sigma & \text{if } \esem{\texttt{E}} \sigma = \supInter{\mathfrak{tt}}\\
        \ssem{\texttt{S}_2} \sigma & \text{if } \esem{\texttt{E}} \sigma = \supInter{\mathfrak{ff}}\\
        undefined & otherwise
      \end{cases}
      &\\[.5em]
    \end{array}&&
    \end{flalign*}
  \end{block}
\end{frame}
%
\begin{frame}{Denotational semantics: loops}
  \begin{itemize}
  \item Loops are trickier...
  \end{itemize}
  %
  \begin{flalign*}
    \hspace*{-.5em}\begin{array}{lll}
      \ssem{\texttt{\textbf{while} (E) \textbf{do} S}} \sigma
      &\triangleq
      &\begin{cases}
        \sigma & \text{if } \esem{\texttt{E}} \sigma = \supInter{\mathfrak{ff}}\\
        (\ssem{\texttt{\textbf{while} (E) \textbf{do} S}} \circ \ssem{\texttt{S}}) \sigma & \text{if } \esem{\texttt{E}} \sigma = \supInter{\mathfrak{tt}}\\
        undefined & otherwise
      \end{cases}
    \end{array}&&
  \end{flalign*}
\end{frame}
%
\begin{frame}[plain]
  \titlepage
\end{frame}
%
\addtocounter{framenumber}{-1}
%
\begin{frame}[allowframebreaks]{References}
  \scriptsize
  \bibliographystyle{apalike}
  \bibliography{../../clip/bibtex/clip/clip.bib,../../clip/bibtex/clip/general.bib,bib.bib}
\end{frame}
%
\end{document}

%%% Local Variables:
%%% mode: LaTeX
%%% TeX-master: t
%%% End:
